\section{Divisão de papéis}

Com três integrantes e seis semanas, o grupo optou por \textbf{papéis
fixos}, um por frente de trabalho, em vez de papéis rotativos:

\begin{table}[H]
  \centering
  \small
  \begin{tabularx}{\textwidth}{>{\raggedright}p{3.6cm}X}
    \toprule
    \textbf{Papel} & \textbf{Responsabilidades} \\
    \midrule
    Dados e documentação &
      Análise e preparo dos datasets; consolidação dos relatórios; slides;
      registro das decisões \\
    Desenvolvimento &
      Código MPI (mestre, workers, fila dinâmica, tolerância a falhas) e
      versão sequencial de referência \\
    Infraestrutura e experimentos &
      Vagrant, Ansible e cluster; execução dos experimentos e coleta de
      métricas \\
    \bottomrule
  \end{tabularx}
  \caption{Papéis fixos do grupo.}
\end{table}

A atribuição de cada papel ainda será definida pelo grupo. Mesmo com
papéis fixos, a tendência é que todos participem um pouco de todas as
frentes: o papel indica quem responde por cada uma, e não um limite para a
colaboração. O trabalho é versionado em um repositório público no
GitHub\footnote{\url{https://github.com/Matheus-Nardi/projeto-sd}}.
